% ===========================================================================
% ΔΟΚΙΜΑΣΤΙΚΟ ΦΥΛΛΟ ΕΡΓΑΣΙΑΣ
% ===========================================================================
% \WorksheetOpening: ξεχωριστή μπλε επικεφαλίδα και αύξων αριθμός φύλλου.
% \WorksheetSection: μπλε μπάρα για επιμέρους δραστηριότητα.
% \CheckSquare: τετραγωνάκι επιλογής, με μέγεθος από book-config.tex.
% \AnswerLine/\AnswerGrid: χώρος για γραπτές απαντήσεις μαθητή.
% ===========================================================================
% The worksheet is intentionally independent of the subject text.
\WorksheetOpening{ΚΙΝΗΜΑΤΙΚΗ}{Θέση, μετατόπιση και ερμηνεία διαγραμμάτων}
\begin{LearningGoals}
\begin{enumerate}
\item Να υπολογίζουν μετατόπιση από δύο θέσεις.
\item Να επιλέγουν σωστή μονάδα φυσικού μεγέθους.
\item Να αιτιολογούν σύντομα τις απαντήσεις τους.
\end{enumerate}
\end{LearningGoals}
\WorksheetSection{Α. Ελέγχω την κατανόησή μου}
% ΣΤΥΛ v0.16:
% Ο τίτλος «Σωστό ή λάθος;» έχει μόνο μαύρα έντονα γράμματα.
% Το υπόλοιπο κείμενο είναι κανονική μαύρη Inter (όχι μπλε ή bold).
% Η μορφοποίηση γίνεται κεντρικά από το exercise στο universalbook.sty.
\begin{exercise}[]{Σωστό ή λάθος;}
Η μετατόπιση ενός κινητού είναι πάντα ίση με το διάστημα που διήνυσε.\hspace{.5em}\hfill\TrueFalseChoices
\end{exercise}
\AnswerLine
\WorksheetSection{Β. Εφαρμόζω}
\begin{exercise}[]{Υπολογισμός μετατόπισης}
Ένα κινητό μετακινείται από τη θέση \(x_A=-2\,\mathrm m\) στη θέση \(x_B=6\,\mathrm m\). Να υπολογίσετε τη μετατόπιση και να σχολιάσετε το πρόσημό της.
\end{exercise}
\AnswerGrid{2}
\WorksheetSection{Γ. Εξηγώ}
\begin{exercise}[]{Ερώτηση κρίσεως}
Να περιγράψετε μια κίνηση στην οποία το διάστημα είναι θετικό, αλλά η συνολική μετατόπιση είναι μηδέν.
\end{exercise}
\AnswerLine
\vfill\begin{center}{\footnotesize\color{BookSecondary}ΔΕΙΓΜΑ ΦΥΛΛΟΥ ΕΡΓΑΣΙΑΣ · ΔΕΝ ΑΠΟΤΕΛΕΙ ΤΕΛΙΚΗ ΔΙΔΑΚΤΙΚΗ ΥΛΗ}\end{center}
